\section{Saran}

Berdasarkan pengalaman selama pelaksanaan magang, terdapat beberapa hal yang
dapat menjadi perhatian bagi mahasiswa yang akan mengikuti program magang pada
posisi yang serupa di masa mendatang:

\begin{enumerate}
      \item Mempelajari dasar-dasar framework dan teknologi yang digunakan oleh perusahaan
            sebelum memulai magang akan sangat membantu dalam mempercepat proses adaptasi.
            Pemahaman mengenai Flutter, React, TypeScript, serta konsep dasar API akan
            memberikan kesempatan yang lebih besar untuk berkontribusi secara aktif sejak
            awal pelaksanaan magang.

      \item Tidak ragu untuk mengambil inisiatif ketika menemukan permasalahan atau peluang
            perbaikan pada sistem yang sedang dikembangkan. Kemampuan untuk
            mengidentifikasi bug, mengusulkan solusi, serta membuat \textit{task} secara
            mandiri merupakan nilai tambah yang sangat dihargai dalam lingkungan kerja
            profesional.

      \item Membiasakan diri untuk mendokumentasikan pekerjaan dan proses penyelesaian
            masalah secara rutin sejak awal pelaksanaan magang. Catatan mengenai solusi
            yang ditemukan, keputusan teknis yang diambil, maupun kendala yang dihadapi
            akan sangat membantu dalam proses evaluasi diri serta penyusunan laporan akhir
            magang.

      \item Memiliki pemahaman mengenai penggunaan Git dan alur kolaborasi pengembangan
            perangkat lunak seperti \textit{branching}, \textit{merge request}, dan
            \textit{code review}. Kemampuan tersebut sangat membantu dalam beradaptasi
            dengan proses pengembangan perangkat lunak yang diterapkan pada lingkungan
            industri.

      \item Menjaga komunikasi yang baik dengan mentor dan anggota tim lainnya. Diskusi
            secara aktif ketika mengalami kesulitan dapat mempercepat proses pembelajaran
            serta membantu memperoleh pemahaman yang lebih baik mengenai kebutuhan bisnis
            dan implementasi teknis yang diterapkan dalam proyek.
\end{enumerate}